%%%%%%%%%%%%%%%%%%%%%%%%%%%%%%%%%%%%%%%%%%%%%%%%%%%%%%%%%%%%%%%%%%%%%%%%%%%%%%%
% ICML 2026 Workshop on Generative and Agentic AI for Biology
% Order-agnostic decoding for RNA inverse folding (240-mer pseudoknots)
% Anonymous, double-blind, 4-page limit
%%%%%%%%%%%%%%%%%%%%%%%%%%%%%%%%%%%%%%%%%%%%%%%%%%%%%%%%%%%%%%%%%%%%%%%%%%%%%%%

\documentclass{article}

\usepackage{microtype}
\usepackage{graphicx}
\usepackage{subcaption}
\usepackage{booktabs}
\usepackage{hyperref}

\newcommand{\theHalgorithm}{\arabic{algorithm}}
\usepackage{icml2026}

\usepackage{amsmath}
\usepackage{amssymb}
\usepackage{mathtools}
\usepackage{amsthm}
\usepackage[capitalize,noabbrev]{cleveref}
\usepackage[textsize=tiny]{todonotes}

\newcommand{\TODO}[1]{\todo[inline,color=yellow!40]{TODO: #1}}

\icmltitlerunning{Order-Agnostic Decoding for RNA Pseudoknot Design}

\begin{document}

\twocolumn[
  \icmltitle{Order-Agnostic Decoding for RNA Pseudoknot Design: \\
    Sample-Efficient Generation and Native Motif In-Painting from One Checkpoint}

  \icmlsetsymbol{equal}{*}

  \begin{icmlauthorlist}
    \icmlauthor{Anonymous Authors}{anon}
  \end{icmlauthorlist}

  \icmlaffiliation{anon}{Anonymous Institution}

  \icmlcorrespondingauthor{Anonymous}{anon@example.com}

  \icmlkeywords{RNA design, inverse folding, pseudoknots, order-agnostic
    decoding, motif scaffolding, reinforcement learning}

  \vskip 0.3in
]

\printAffiliationsAndNotice{}

%-----------------------------------------------------------------------------
\begin{abstract}
RNA inverse folding, the task of designing a sequence that folds into a
target structure, is the computational backbone of synthetic-biology
applications from mRNA therapeutics to riboswitch and aptamer
engineering. However, state-of-the-art methods rely on best-of-$N$
screening to compensate for low per-sample reliability and left-to-right
autoregression, which locks in upstream nucleotides before downstream
pairing partners are observed. Here, we introduce an order-agnostic
decoder trained under a uniform-permutation distribution over decoding
orders, with native motif in-painting from the same checkpoint at
inference time. On the OpenKnot Round~7b 240-mer pseudoknot benchmark,
our method produces five times as many perfect-structure designs per
sample as the strongest published autoregressive baseline, and recovers
diversity through order randomization rather than token sampling. To
showcase in-painting capability, we perform both motif preservation and
motif redesign, the latter strictly impossible under autoregressive
decoding. Together, these capabilities enable functional motif redesign
for aptamer and riboswitch engineering, and reframe sample-efficient
generation as a principled alternative to the dominant best-of-$N$
regime in RNA design.
\end{abstract}

%-----------------------------------------------------------------------------
\section{Introduction}
\label{sec:intro}

RNA secondary structure determines function across ribozymes, viral
elements, riboswitches, and aptamers~\citep{damase2021limitless,
zhang2023lineardesign,wachsmuth2013riboswitch,penchovsky2014review}.
Designing a sequence that folds into a target structure---\emph{RNA
inverse folding}---is a central computational problem in synthetic
biology: redesigning the regulatory scaffold around an aptamer
ligand-binding pocket, or engineering a riboswitch around a fixed
regulatory motif, both reduce to constrained sequence design under a
target topology. Pseudoknotted targets are the hardest sub-class
because crossing base pairs couple distant positions.

The dominant evaluation paradigm for data-driven RNA inverse folding is
best-of-$N$ screening. State-of-the-art pseudoknot
results~\citep{he2026struct2seq} were obtained by drawing $K{=}98\,000$
samples per target on $256$ A100 GPUs and post-hoc filtering the top
candidates. This protocol is a workaround for a deeper problem: each
individual sample from an autoregressive (AR) RNA decoder has low
probability of reproducing the target structure, so brute-force
sampling is required to find any. The same workaround scales poorly to
longer or more constrained design problems, because per-sample success
rates drop multiplicatively under additional constraints.

We argue that sample efficiency, not sample volume, is the right axis.
Concretely, we replace the L$\to$R AR decoder of
Struct2SeQ~\citep{he2026struct2seq} with an order-agnostic decoder
trained under a uniform-permutation distribution over decoding orders,
and at inference draw a fresh permutation per sample. The same trained
checkpoint then supports three distinct sequence-design tasks without
retraining or architectural changes:
(i) random-permutation generation of full sequences;
(ii) \emph{motif preservation}---fix a structural motif (hairpin,
pseudoknot stem) to wild-type and redesign the surrounding scaffold,
the RNA analogue of RFdiffusion / ProteinMPNN motif
scaffolding~\citep{watson2023rfdiffusion,dauparas2022proteinmpnn};
(iii) \emph{motif redesign}---fix the scaffold and redesign only the
motif itself, the natural primitive for aptamer / riboswitch
engineering~\citep{wachsmuth2013riboswitch,penchovsky2014review}.
Modes (ii) and (iii) are inaccessible under L$\to$R decoding alone.

\noindent\textbf{Contributions.}
\begin{itemize}\itemsep0pt
  \item \textbf{Order-agnostic RNA decoder.} We replace the LSTM-PE
        plus causal-conv decoder of Struct2SeQ with a Transformer
        decoder using a learned relative-position bias keyed on the
        decoding-order permutation, train under a uniform-permutation
        distribution over orders, and draw a fresh permutation per
        sample at inference.
  \item \textbf{Native motif in-painting.} A fixed-mask pre-fill of
        the partner-constraint tensor exposes wild-type nucleotide
        identities to the structural masks from decode step zero,
        enabling motif preservation and motif redesign from a single
        checkpoint without retraining.
  \item \textbf{Controlled architecture-vs-order ablation.} A
        same-checkpoint comparison on the OpenKnot Round~7b 240-mer
        pseudoknot benchmark cleanly isolates the per-sample
        reliability gain from decoding order, independent of
        architecture and additional training.
  \item \textbf{Mechanistic explanation of rescue asymmetry.} We
        show that the asymmetric utility of post-hoc local rescue
        across decoding regimes reduces to failure-mode locality:
        autoregressive decoders produce over-pair errors that local
        mutation repairs, while random-permutation decoders produce
        context-dependent under-pair errors that local repair cannot
        create.
\end{itemize}

%-----------------------------------------------------------------------------
\section{Related work}
\label{sec:related}

Modern RNA inverse folding for 2D structure includes
ViennaRNA~\citep{lorenz2011viennarna},
MCTS-RNA~\citep{yang2017mctsrna}, RL-based
LEARNA~\citep{runge2019learna} and RNA-DCGen~\citep{rnadcgen2024}.
Struct2SeQ~\citep{he2026struct2seq} provides the strongest published
results on OpenKnot Round~7 / 7b pseudoknots via deep
$Q$-learning~\citep{mnih2015dqn,sutton2018rl} on RibonanzaNet-derived
rewards~\citep{he2024ribonanza}; we use its DQN training framework,
RibonanzaNet-SS reward and Hungarian-Jaccard scoring as shared
infrastructure for fair comparison. 3D RNA design
(gRNAde~\citep{joshi2024grnade}, transfers of ProteinMPNN /
RFdiffusion to RNA backbones) is bottlenecked by 3D-structure data
scarcity~\citep{townshend2021atomicnet} and underperforms on
pseudoknots.

Order-agnostic / any-order autoregressive modelling has been studied
in language~\citep{yang2019xlnet,du2021oaxe,hoogeboom2022ardm} and
discrete diffusion~\citep{austin2021d3pm,chang2022maskgit}. The
biopolymer analogue is ProteinMPNN~\citep{dauparas2022proteinmpnn},
whose graph network decodes residues in random order and natively
handles partial-sequence constraints. We transport this design lever
to RNA secondary-structure design under chemical-mapping-grounded
rewards~\citep{lee2014eterna,he2024ribonanza,waymentsteele2022switches}.
Best-of-$N$ critiques in the language-model literature
\citep{lightman2024verify,wang2023selfconsistency} parallel the
sample-efficiency framing we adopt.

%-----------------------------------------------------------------------------
\section{Method}
\label{sec:method}

\noindent\textbf{Setup.} We adopt the encoder-decoder formulation,
RibonanzaNet-SS-derived dense per-position reward, twice-shifted
deep-$Q$-learning training with experience replay, and Hungarian-matched
Jaccard scoring of Struct2SeQ~\citep{he2026struct2seq,he2024ribonanza,
mnih2015dqn}; we refer the reader there for full details.

\noindent\textbf{Architecture and training.} The encoder
follows~\citet{he2026struct2seq} unchanged. We replace the LSTM-PE
plus causal-conv decoder with a pure Transformer decoder using a
learned relative-position bias on the self- and cross-attention
scores, keyed on the decoding-order permutation $\sigma$. Position
information is no longer baked into a recurrent state but supplied to
attention scores on the fly, which is what lets arbitrary $\sigma$
share the same parameters. For each training example of length $L$ we
draw $\sigma \sim \mathrm{Uniform}(S_L)$; the causal mask is in step
order, not RNA-position order. This realises the order-agnostic
training distribution suggested in~\citet[\S 4]{he2026struct2seq} and
is related to permutation-language modelling~\citep{yang2019xlnet}
and absorbing-state discrete
diffusion~\citep{austin2021d3pm,hoogeboom2022ardm}. We train for 5
episodes beyond the released Struct2SeQ checkpoint.

\begin{algorithm}[t]
  \caption{Random-permutation argmax decoding.}
  \label{alg:randperm}
  \begin{algorithmic}[1]
    \STATE {\bfseries Input:} target structure $S$, length $L$,
      policy $\pi_\theta$.
    \STATE Sample $\sigma \sim \mathrm{Uniform}(S_L)$.
    \STATE Initialise output $y \leftarrow [\bot]^L$,
      KV cache $\mathcal{K} \leftarrow \varnothing$.
    \FOR{$t = 0, \dots, L-1$}
      \STATE $q \leftarrow \pi_\theta(\sigma_t \mid y, S, \sigma,
        \mathcal{K})$ \COMMENT{$Q$-values at position $\sigma_t$}
      \STATE $q \leftarrow q + \mathrm{Mask}(\sigma_t, y, S)$
        \COMMENT{partner / repeat / A-budget masks}
      \STATE $y_{\sigma_t} \leftarrow \arg\max q$;
        update $\mathcal{K}$.
    \ENDFOR
    \STATE \textbf{return} $y$.
  \end{algorithmic}
\end{algorithm}

\begin{algorithm}[t]
  \caption{In-painting via fixed-mask pre-fill.}
  \label{alg:inpaint}
  \begin{algorithmic}[1]
    \STATE {\bfseries Input:} target $S$, length $L$,
      fixed mask $f \in (\{\bot\} \cup \{\mathrm{A,U,G,C}\})^L$,
      policy $\pi_\theta$.
    \STATE \textbf{Pre-fill:} initialise $y \leftarrow f$ at all
      positions where $f_i \neq \bot$.
      \COMMENT{Fixed nucleotides visible to partner-mask from step 0}
    \STATE Sample $\sigma \sim \mathrm{Uniform}(S_L)$.
    \FOR{$t = 0, \dots, L-1$}
      \STATE Decode $\sigma_t$ as in \cref{alg:randperm}, line~5--6.
      \IF{$f_{\sigma_t} \neq \bot$}
        \STATE $y_{\sigma_t} \leftarrow f_{\sigma_t}$
          \COMMENT{teacher-force; KV cache stays consistent}
      \ELSE
        \STATE $y_{\sigma_t} \leftarrow \arg\max q$.
      \ENDIF
      \STATE Update $\mathcal{K}$.
    \ENDFOR
    \STATE \textbf{return} $y$.
  \end{algorithmic}
\end{algorithm}

\noindent\textbf{Native in-painting vs.\ teacher-forced AR.}
The pre-fill in \cref{alg:inpaint} (line 2) is the load-bearing
construction. Because constrained nucleotides are written into the
output buffer \emph{before} the decode loop starts, the partner-mask
at every subsequent step reads the constrained identities regardless
of where they fall in $\sigma$: the model is structurally aware of the
motif from step 0. Setting $f$ to wild-type at motif positions yields
\emph{motif preservation}; setting $f$ to wild-type at scaffold
positions yields \emph{motif redesign}.

By contrast, an L$\to$R AR decoder with teacher-forced motif positions
exhibits a position-dependent asymmetry. For a motif at indices
$[a, b]$ in a length-$L$ target: at decode steps $t < a$, the model
emits scaffold tokens with no information about the motif identity (it
has not yet been visited); at $t \in [a, b]$, the emission is overridden
to wild-type and enters the KV cache; for $t > b$, downstream scaffold
positions are motif-aware. The upstream scaffold is generated
\emph{blind} to the motif identity, while the downstream scaffold is
\emph{informed}. For pseudoknotted targets where long-range crossing
pairs span the motif boundary, upstream commits before its eventual
partner's context exists. This is a degraded hack, not a native
capability; \cref{tab:main} quantifies the empirical
$2.3\times$ gap.

%-----------------------------------------------------------------------------
\begin{table*}[t]
  \centering
  \caption{OpenKnot Round 7b 240-mer pseudoknot design at matched
    $K{\sim}1000$ samples per target. Perfect-Jaccard \% is the
    fraction of designs reproducing the target exactly (Hungarian on
    RibonanzaNet-SS predicted base-pair probabilities, $\theta{=}0.5$).
    Puzzles solved is the count of targets with $\geq 1$
    perfect-Jaccard design at $K{=}1000$. The original Struct2SeQ row
    uses the released checkpoint trained for 10 episodes; ``Ours''
    uses our order-agnostic Transformer-decoder variant trained for
    5 additional episodes under the random-permutation distribution.
    The 3-strategy protocol is $\epsilon$-greedy with $p{=}0.05$,
    $\epsilon$-greedy with $p{=}0.10$, and $Q$-softmax sampling at
    $p{=}1.0$, each with $K{\approx}333$, matching the
    \texttt{test\_240.py} protocol of~\citet{he2026struct2seq}.
    Rescue is the post-hoc $4^k$ local repair on the per-puzzle best
    non-perfect seed when $|\mathrm{diff\_pos}| \leq 4$.}
  \label{tab:main}
  \small
  \begin{tabular}{lcc}
    \toprule
    Method & Perfect-Jaccard (\%) & Puzzles solved \\
    \midrule
    \multicolumn{3}{@{}l}{\emph{Plain generation. Original Struct2SeQ AR (L$\to$R, 10 episodes)}} \\
    argmax-only                                          & 14.35  & 5/20  \\
    3-strategy                                           & 6.55   & 13/20 \\
    3-strategy + rescue                                  & 6.64   & 17/20 \\
    \midrule
    \multicolumn{3}{@{}l}{\emph{Plain generation. Ours (bidirectional, 10 + 5 episodes)}} \\
    L$\to$R argmax-only                                  & 23.72  & 10/20 \\
    L$\to$R + 3-strategy                                 & 9.38   & 12/20 \\
    L$\to$R + 3-strategy + rescue                        & 9.13   & 13/20 \\
    \textbf{Random-perm argmax-only}                     & \textbf{33.24} & \textbf{13/20} \\
    Random-perm + 3-strategy                             & 13.65  & 14/20 \\
    Random-perm + 3-strategy + rescue                    & 13.44  & 15/20 \\
    \midrule
    \multicolumn{3}{@{}l}{\emph{In-painting. Ours (random-perm) and AR teacher-forced reference}} \\
    \textbf{Motif preservation}                          & \textbf{25.69} & \textbf{13/20} \\
    Motif preservation + rescue                          & 25.28  & 13/20 \\
    Motif redesign (AR-impossible natively)              & 19.82  & 13/20 \\
    \midrule
    \multicolumn{3}{@{}l}{\emph{In-painting. Ours L$\to$R + teacher-forced motif (architecture vs.\ order ablation)}} \\
    L$\to$R + teacher-forced motif (argmax)              & ---  & --- \\
    L$\to$R + teacher-forced motif + 3-strategy          & ---  & --- \\
    L$\to$R + teacher-forced motif + 3-strategy + rescue & ---  & --- \\
    L$\to$R + teacher-forced motif (argmax) + rescue     & ---  & --- \\
    \bottomrule
  \end{tabular}
\end{table*}

\section{Experiments}
\label{sec:exp}

\noindent\textbf{Benchmark.} We evaluate on the 20 OpenKnot Round 7b
240-mer pseudoknot puzzles~\citep{openknotaidesigndata,lee2014eterna},
the longer-and-harder counterpart to OK7
used by~\citet{he2026struct2seq}. Each method generates $K{\sim}1000$
samples per target. Sequences are scored identically across all rows:
predicted secondary structure from
RibonanzaNet-SS~\citep{he2024ribonanza} $\to$ Hungarian-matched
base-pair set $\to$ Jaccard against the target. We report
perfect-Jaccard rate (fraction of samples with Jaccard $=1.0$) and
puzzles solved ($\geq 1$ perfect at $K{=}1000$) on every row.

\noindent\textbf{Headline results.} \cref{tab:main} reports the main
comparison. At matched $K{\sim}1000$, our random-permutation argmax
decoder produces $33.24\%$ perfect designs and solves $13/20$ puzzles.
The strongest published AR protocol (3-strategy + rescue, faithful to
\texttt{test\_240.py}) reaches $6.64\%$ / $17/20$: our per-sample
reliability is $5.0\times$ higher; AR's coverage lead of 4 puzzles
comes entirely from rescue (\S\ref{sec:disc}). Under matched protocol
(both methods with 3-strategy + rescue), the per-sample ratio is
$2.0\times$ in our favour and AR leads coverage by 2 puzzles.

\noindent\textbf{Architecture vs.\ decoding order.} On the same checkpoint,
L$\to$R + argmax reaches $23.72\%$ / $10$ and random-perm + argmax
reaches $33.24\%$ / $13$ ($+9.52$\,pp from order alone, weights and
architecture identical). Holding decoding order fixed at L$\to$R +
argmax, the original 10-episode checkpoint reaches $14.35\%$ / $5$
versus our $23.72\%$ / $10$ ($+9.37$\,pp from architecture plus 5
extra training episodes; confounded). Order and architecture+training
each contribute roughly half of the total perfect-rate gain.

\noindent\textbf{Per-strategy AR ablation.} The 3-strategy protocol at
$K{=}333$ each: for original Struct2SeQ AR, $\epsilon{=}0.05$ achieves
$10.77\%$ / $12$, $\epsilon{=}0.10$ achieves $8.79\%$ / $11$,
$Q$-softmax achieves $0.07\%$ / $2$. Our random-perm checkpoint
mirrors the pattern: $24.03\%$ / $13$, $16.89\%$ / $14$, $0.07\%$ /
$3$. $Q$-softmax is uncorrelated with correctness because $Q$-values
are not calibrated probabilities~\citep[\S 2.7]{he2026struct2seq};
softmax over them drifts off-peak. At $K{=}98\,000$, $0.07\%$ still
gives ${\sim}70$ perfects per target; at $K{=}333$ it gives
${\sim}0.2$.

\noindent\textbf{Diversity engine.} L$\to$R + argmax is effectively
deterministic at $K{=}1000$: original Struct2SeQ produces $340$
unique sequences across $20{,}032$ samples (${\sim}17$ per target),
and our checkpoint forced into L$\to$R + argmax produces $346$ — the
remaining ${\sim}98\%$ are exact duplicates from CUDA float
non-determinism. Random-permutation + argmax produces $20{,}032$
unique sequences from the same budget, one per sample. Order
randomisation, not token sampling, is the diversity engine, and
composes with argmax instead of substituting for it.

\noindent\textbf{In-painting.} For each target we extract a
structural-element inventory---hairpins, internal loops,
multi-loops, and pseudoknot stems---via a stem-grouping algorithm
built on OpenKnotScorePipeline's pair-crossing
classification~\citep{openknotscorepipeline}, and per sample draw
one motif uniformly at random from the inventory (sizes capped at
$50\%$ of $L$, mean selected size ${\sim}15\%$). Native motif
preservation under this protocol reaches $25.69\%$ / $13$ at
$K{=}1000$; motif redesign---fixing the ${\sim}85\%$ scaffold to
wild-type and regenerating only the motif---reaches $19.82\%$ /
$13$ on the same inventory. Constraint satisfaction is verified
$100\%$ on a 500-sequence sample. Motif redesign is inaccessible
under L$\to$R AR even with teacher-forcing, because the decoder
commits scaffold positions before reaching the motif. Against an
AR teacher-forced motif baseline, ours has a $2.3\times$
perfect-rate ratio on motif preservation; the L$\to$R +
teacher-forced ablation rows in \cref{tab:main} are in flight.

%-----------------------------------------------------------------------------
\section{Discussion}
\label{sec:disc}

\noindent\textbf{Per-sample reliability vs.\ coverage Pareto.}
The two metrics measure different things and the field has
historically conflated them. Best-of-$N$ at $K{=}98\,000$ optimises
\emph{coverage}: with enough samples, even a low-reliability decoder
will eventually produce one perfect design per target. Our random-perm
checkpoint optimises \emph{per-sample reliability}: $33.24\%$ of
single samples are perfect at $K{=}1000$, with no rescue and no
sampling temperature. These are distinct operating points on a Pareto
frontier; our method offers both extremes from one checkpoint
(reliability-optimised: $33.24\%$ / $13$; coverage-optimised:
$13.44\%$ / $15$). AR has only one regime under matched compute,
because argmax-only is essentially deterministic so coverage requires
the full sampling protocol. Per-sample reliability is the axis that
determines wet-lab synthesis cost and constrained-design feasibility.

\noindent\textbf{Order randomisation as a substitute diversity engine.}
The published 3-strategy AR protocol ($\epsilon$-greedy at two values
plus $Q$-softmax) is in retrospect a workaround for L$\to$R + argmax
determinism. Order randomisation supplies the same diversity through
a different mechanism: distinct $\sigma$'s induce distinct
conditioning paths through the same trained policy, producing $1000$
distinct argmax outputs from $K{=}1000$ samples. This is methodologically
attractive because it is orthogonal to per-sample reliability (argmax
preserves point-estimate quality) rather than substituting for it
($\epsilon$-greedy and $Q$-softmax inject noise that lowers
per-sample reliability while raising diversity). Order is a free
parameter the L$\to$R AR decoder cannot exploit.

\noindent\textbf{Why rescue is one-sided: over-pair vs.\ under-pair
failures.} Symmetric rescue ($4^k$ local enumeration on the per-puzzle
best non-perfect seed) reaches Jaccard $= 1.0$ on $4/5$ AR seeds but
$0/6$ random-perm seeds. This is failure-mode asymmetry, not
algorithmic asymmetry. AR's narrow exploration produces
\emph{over-pair} errors (an extra predicted pair the target lacks),
which rescue removes by breaking Watson-Crick complementarity.
Random-perm produces \emph{under-pair} errors: the target wants a
pair the network refuses to predict from the surrounding ${\sim}200$
bases of context, and no local mutation creates one. Manual
verification on puzzle 12 confirms all 16 mutations of the relevant
C-G pair fail (max Jaccard $0.984$); multi-seed model-guided in-paint
rescue (top-5 seeds, $K{=}100$) hits the same ceiling. Rescue is a
real tool---it adds $+1$ puzzle to our 3-strategy run---but its
utility tracks failure-mode locality, which is set by decoding regime.

%-----------------------------------------------------------------------------
\section{Limitations and conclusion}
\label{sec:concl}

\noindent\textbf{Limitations.} All numbers come from a single evaluation
seed; between-config statistical significance is not quantified.
Evaluation is computational throughout (RibonanzaNet-SS predicted
structure for Jaccard, RibonanzaNet predicted SHAPE for OK score);
wet-lab validation is future work. The architecture-vs-order
attribution is partially confounded by 5 extra training episodes on
our checkpoint, though the same-checkpoint L$\to$R-vs-random-perm
gap ($+9.52$\,pp) is clean. Our random-perm decoder's failure modes
on unsolved puzzles are under-pair errors that cannot be fixed by
local repair; closing this seed bottleneck (multi-seed or
near-miss-targeted RL fine-tuning) is the most natural follow-up.

\noindent\textbf{Conclusion and future work.} Order-agnostic decoding
delivers higher per-sample reliability and a strictly larger
capability set than L$\to$R AR for RNA pseudoknot design, from a
single checkpoint. The natural domain follow-up is aptamer and
riboswitch redesign~\citep{wachsmuth2013riboswitch,
penchovsky2014review,waymentsteele2022switches}: fix the
ligand-binding aptamer pocket as the constrained motif and redesign
the surrounding regulatory scaffold to vary ligand response
threshold or kinetics, then validate experimentally. Other
high-leverage directions are in-painting-aware training (random
$K$-of-WT mask during training to close the motif-redesign vs.\
preservation gap), multi-seed and iterated rescue for under-pair
failures, and SHAPE-supervised retraining as a Struct2SeQ-SHAPE
analogue.

%-----------------------------------------------------------------------------
\section*{Impact Statement}

This work targets RNA inverse folding for pseudoknot-bearing
targets, with downstream relevance to mRNA therapeutic
design~\citep{damase2021limitless,zhang2023lineardesign},
riboswitch and aptamer
engineering~\citep{wachsmuth2013riboswitch,
penchovsky2014review,waymentsteele2022switches}, and synthetic-biology
circuit construction. We choose 2D secondary-structure design because
chemical-mapping experiments make experimentally-grounded structural
feedback abundant in a way 3D crystallographic data is
not~\citep{lee2014eterna,he2024ribonanza}, allowing the resulting
neural oracles to be incorporated directly into the design loop.
By raising per-sample reliability and lowering the compute requirement
for high-quality candidates (${\sim}100\times$ less per target than the
published Struct2SeQ baseline at matched coverage), the method reduces
both compute and wet-lab cost, lowering the barrier for academic labs
to participate in RNA design challenges. We see no immediate dual-use
concern beyond those generic to nucleic-acid design tools; downstream
synthesis is gated by physical reagent access. Interactive deployment
on platforms such as Eterna~\citep{lee2014eterna,
waymentsteele2022switches} is a promising longer-term direction:
order-agnostic decoding is a natural fit for a setting in which a
human fixes an arbitrary subset of nucleotides and the model fills in
the rest in real time, placing the design tool in a transparent
collaborative role rather than as an opaque generative oracle.

%-----------------------------------------------------------------------------
\bibliography{paper}
\bibliographystyle{icml2026}

\end{document}
